\chapter{Complexity — ZPP Hosts Classes}
%%% generated by the blueprint pipeline from TCSlib.Complexity.Randomized.ZPPHostsClasses
\section{Overview}

Declarations of the Lean module \texttt{TCSlib.Complexity.Randomized.ZPPHostsClasses}.

\section{Declarations}

\begin{theorem}[Power bound for a shifted natural number]
\lean{Randomized.succ_pow_le_two_pow_mul}
\label{Randomized.succ_pow_le_two_pow_mul}
\leanok
\difficulty{4}
For all natural numbers $n$ and $p$, $(n+1)^p \le 2^p(n^p+1)$.
\end{theorem}

\begin{theorem}[Polynomial power inequality]
\lean{Randomized.add_mul_succ_pow_le}
\label{Randomized.add_mul_succ_pow_le}
\leanok
\uses{Randomized.succ_pow_le_two_pow_mul}
\difficulty{4}
For all natural numbers $A,B,C,n,p$, $$A(n+1)^{p+1}+B(n+1)^p+C\le
(A+B+C)2^{p+1}\bigl(n^{p+1}+1\bigr).$$
\end{theorem}

\begin{theorem}[Polynomial-time aborting decision procedure]
\lean{Turing.FinNDTM.ZDecidesInTime.exists_abortDecidesInTime}
\label{Turing.FinNDTM.ZDecidesInTime.exists_abortDecidesInTime}
\leanok
\uses{Randomized.add_mul_succ_pow_le, Randomized.attemptBit_eq_some, Randomized.clockHost_abortDecides, Randomized.randProb_congr, Turing.FinNDTM, Turing.FinNDTM.AbortDecidesInTime, Turing.FinNDTM.AbortDecidesInTime.mono, Turing.FinNDTM.ExpTimeLe.one_le, Turing.FinNDTM.ZDecidesInTime, Turing.FinNDTM.ZDecidesInTime.truncate, Turing.FinNDTM.haltRun_eq_some_iff, Turing.FinNDTM.outputWith, Turing.FinTM.ComputesFunInTime, Turing.FinTM.computesFunInTime_polyUnary}
\difficulty{6}
Let $N$ be a finite nondeterministic Turing machine over Boolean symbols, let $L$ be a
language of Boolean words, and let $c,e\in\mathbb{N}$. Suppose that on every input $x$,
the sum of the probabilities that the run has not halted, over every finite range of
time horizons, is at most $c(|x|^e+1)$, and every halted run outputs the one-symbol
indicator of whether $x\in L$. Then there exist $K,E\in\mathbb{N}$ and a finite
nondeterministic Turing machine $N'$ over Boolean symbols such that, on every input $x$,
$N'$ halts on every branch after $K(|x|^E+1)$ steps, its output on every branch of that
length differs from the one-symbol indicator of membership in $L$, and the probability
that its output equals that indicator is at least $1/2$.
\end{theorem}

\begin{theorem}[Zero-error classes with and without aborts]
\lean{Randomized.ZPP_PTM_subset_ZPPAbort_PTM}
\label{Randomized.ZPP_PTM_subset_ZPPAbort_PTM}
\leanok
\uses{Randomized.ZPPAbort_PTM, Randomized.ZPP_PTM, Turing.FinNDTM.ZDecidesInTime.exists_abortDecidesInTime}
\difficulty{3}
Every language decidable in zero-error expected polynomial time belongs to the
polynomial-time abort class.
\end{theorem}

\begin{definition}[Conditional Boolean output]
\lean{Randomized.keepAs}
\label{Randomized.keepAs}
\leanok
For a Boolean $b$, define a function on Boolean inputs $eta$ by returning the optional
value $b$ when $eta$ is true, and no value when $eta$ is false.
\end{definition}

\begin{definition}[Zero-error randomized attempt machine]
\lean{Randomized.zppAttempt}
\label{Randomized.zppAttempt}
\leanok
\uses{Randomized.keepAs, Turing.FinNDTM, Turing.FinNDTM.coinBranch, Turing.FinNDTM.mapOutput}
Given a function $N$ assigning a finite nondeterministic Turing machine over the Boolean
alphabet to each Boolean choice, form the machine obtained by mapping the outputs of
$N(	ext{false})$ and $N(	ext{true})$ through the conditional Boolean output functions
for false and true, respectively, and combining the two machines by the coin-branch
construction.
\end{definition}

\begin{theorem}[Halting of the selected machine branch]
\lean{Randomized.zppAttempt_halted}
\label{Randomized.zppAttempt_halted}
\leanok
\uses{Randomized.keepAs, Randomized.zppAttempt, Turing.FinNDTM, Turing.FinNDTM.coinBranchCfg₀, Turing.FinNDTM.coinBranchCfg₁, Turing.FinNDTM.coinBranch_runWith, Turing.NDTM.initCfg, Turing.NDTM.mapOutput, Turing.NDTM.mapOutputCfg, Turing.NDTM.mapOutput_initCfg, Turing.NDTM.mapOutput_runWith, Turing.NDTM.runWith}
\difficulty{4}
Let $N$ assign a finite nondeterministic Turing machine to each bit, and let $x$ and $u$
be finite Boolean strings and $b$ a bit. The machine formed by the randomized branch
construction halts on input $x$ after the choice string consisting of $b$ followed by
$u$ if and only if the machine $N(b)$ halts on input $x$ after the choice string $u$.
\end{theorem}

\begin{theorem}[Output of a randomized machine attempt]
\lean{Randomized.zppAttempt_output}
\label{Randomized.zppAttempt_output}
\leanok
\uses{Randomized.keepAs, Randomized.zppAttempt, Turing.FinNDTM, Turing.FinNDTM.coinBranchCfg₀, Turing.FinNDTM.coinBranchCfg₁, Turing.FinNDTM.coinBranch_runWith, Turing.NDTM.initCfg, Turing.NDTM.mapOutput, Turing.NDTM.mapOutput_initCfg, Turing.NDTM.mapOutput_runWith, Turing.NDTM.runWith}
\difficulty{4}
Let $N$ assign a finite nondeterministic Turing machine to each bit, let $x$ be a
Boolean input, let $b$ be a Boolean choice, and let $u$ be a finite word of Boolean
choices. The output produced by the randomized attempt on input $x$ with choice word
$b::u$ equals the output of machine $N(b)$ on $x$ with choice word $u$, after retaining
each true output symbol and replacing it by $b$, and discarding each false output
symbol.
\end{theorem}

\begin{theorem}[Zero-error algorithms in expected polynomial time]
\lean{Randomized.inZPP_subset_ZPP_PTM}
\label{Randomized.inZPP_subset_ZPP_PTM}
\leanok
\uses{Randomized.InZPP, Randomized.ZPP_PTM, Randomized.add_mul_succ_pow_le, Randomized.attemptBit, Randomized.attemptBit_eq_some, Randomized.clockHost_zDecides, Randomized.exists_FinNDTM_of_eff, Randomized.keepAs, Randomized.polyLen, Randomized.polyTimeModel, Randomized.polyTimeModel_closedUnderAnswerIs, Randomized.randProb, Randomized.randProb_append_ge, Randomized.randProb_congr, Randomized.randProb_mono, Randomized.randProb_not, Randomized.randProb_one_eq, Randomized.zppAttempt, Randomized.zppAttempt_halted, Randomized.zppAttempt_output, Turing.FinNDTM, Turing.FinNDTM.ZDecidesInTime.mono, Turing.FinNDTM.outputProb, Turing.FinNDTM.outputProb_eq_of_haltsWithin, Turing.FinNDTM.outputWith, Turing.FinNDTM.runWith_eq_runWith_take, Turing.FinTM.ComputesFunInTime, Turing.FinTM.computesFunInTime_polyUnary, Turing.MultiTapeTM.indicator, Turing.NDTM.HaltsWithin, Turing.NDTM.initCfg, Turing.NDTM.runWith}
\difficulty{7}
Every language of bit strings that has an efficient polynomial-witness zero-error
verifier belongs to the class of languages decidable by a finite nondeterministic Turing
machine with expected running time bounded by a polynomial.
\end{theorem}

\begin{theorem}[Verifier characterization of zero-error polynomial time]
\lean{Randomized.ZPP_PTM_eq_setOf_inZPP}
\label{Randomized.ZPP_PTM_eq_setOf_inZPP}
\leanok
\uses{Randomized.InZPP, Randomized.ZPP_PTM, Randomized.inZPP_of_mem_ZPP_PTM, Randomized.inZPP_subset_ZPP_PTM, Randomized.polyTimeModel}
\difficulty{1}
A language over binary strings belongs to zero-error expected polynomial time exactly
when there are a polynomial-time verifier and natural numbers $a,k$ such that, on every
input $x$, the verifier never returns false if $x$ is in the language and never returns
true if $x$ is outside it, for every random string of length $a(|x|+1)^k$, and returns
no answer with probability at most $1/2$ under a uniformly random string of that length.
Here polynomial time for the verifier means that its true and false outputs are each
recognized by a language in polynomial time using the specified self-delimiting pairing
of input and random string.
\end{theorem}

\begin{theorem}[Characterization of ZPP with Polynomial-Time Verifiers]
\lean{Randomized.mem_ZPP_PTM_iff}
\label{Randomized.mem_ZPP_PTM_iff}
\leanok
\uses{Randomized.InZPP, Randomized.ZPP_PTM, Randomized.ZPP_PTM_eq_setOf_inZPP, Randomized.polyTimeModel}
\difficulty{1}
A language of bit strings belongs to the class defined by polynomial-time zero-error
probabilistic machines if and only if there is a polynomial-time verifier and natural
numbers $a,k$ such that, on each input $x$, the verifier uses $a(|x|+1)^k$ random bits,
never outputs false on any random string when $x$ is in the language, and never outputs
true on any random string when $x$ is outside it. On every input, the fraction of random
strings for which it does not halt is at most $1/2$; polynomial-time efficiency means
that the sets of encoded input-randomness pairs giving its true and false outputs both
lie in $P$.
\end{theorem}

\begin{theorem}[Equality of two zero-error classes]
\lean{Randomized.ZPPAbort_PTM_eq_ZPP_PTM}
\label{Randomized.ZPPAbort_PTM_eq_ZPP_PTM}
\leanok
\uses{Randomized.ZPPAbort_PTM, Randomized.ZPP_PTM, Randomized.ZPP_PTM_subset_ZPPAbort_PTM, Randomized.inZPP_of_mem_ZPPAbort_PTM, Randomized.inZPP_subset_ZPP_PTM}
\difficulty{1}
The class of languages decidable with polynomial time and the stated abort allowance
equals the class decidable in zero-error expected polynomial time.
\end{theorem}

\begin{theorem}[Zero-error algorithms with bounded abort time]
\lean{Randomized.inZPP_subset_ZPPAbort_PTM}
\label{Randomized.inZPP_subset_ZPPAbort_PTM}
\leanok
\uses{Randomized.InZPP, Randomized.ZPPAbort_PTM, Randomized.ZPP_PTM_subset_ZPPAbort_PTM, Randomized.inZPP_subset_ZPP_PTM, Randomized.polyTimeModel}
\difficulty{1}
Every language of bit strings that has an efficient polynomial-witness zero-error
verifier belongs to the class of languages decided by an aborting probabilistic Turing
machine with a polynomial time bound, with correctness probability at least $1/2$.
\end{theorem}

\begin{theorem}[Characterization of zero-error polynomial time]
\lean{Randomized.ZPP_PTM_eq_RP_PTM_inter_coRP_PTM}
\label{Randomized.ZPP_PTM_eq_RP_PTM_inter_coRP_PTM}
\leanok
\uses{Randomized.RP_PTM, Randomized.ZPP_PTM, Randomized.ZPP_PTM_subset_RP_PTM_inter_coRP_PTM, Randomized.coRP_PTM, Randomized.inZPP_subset_ZPP_PTM, Randomized.mem_RP_PTM_iff, Randomized.mem_coRP_PTM_iff, Randomized.zpp_eq_rp_inter_corp_polyTime}
\difficulty{1}
The class of languages decidable in zero-error expected polynomial time equals the
intersection of RP and coRP.
\end{theorem}
